\documentclass[11pt]{article}
% Shared print layout. The web's provisional design and MathJax are independent.
\usepackage[utf8]{inputenc}
\usepackage[T1]{fontenc}
\usepackage[spanish,es-nodecimaldot]{babel}
\usepackage{lmodern}
\usepackage{amsmath,amssymb}
\usepackage{graphicx}
\usepackage{geometry}
\geometry{a4paper,left=20mm,right=20mm,top=25mm,bottom=22mm,headheight=16pt,headsep=8mm}
\usepackage{fancyhdr,lastpage}
\usepackage{enumitem}
\usepackage{xcolor}
\usepackage{pgfplots}
\pgfplotsset{compat=1.18}
\usepgfplotslibrary{fillbetween}
\usepackage[hidelinks,pdfencoding=auto]{hyperref}
\hypersetup{pdflang=es-ES,pdfauthor={},pdfsubject={Matemáticas Aplicadas a las Ciencias Sociales II}}
\pagestyle{fancy}
\fancyhf{}
\fancyhead[L]{\small Matemáticas Aplicadas a las Ciencias Sociales II}
\fancyhead[R]{\small 2.\textsuperscript{o} Bachillerato}
\fancyfoot[L]{\footnotesize Matrices III | Edición revisada 1.0}
\fancyfoot[R]{\footnotesize Página \thepage\ de \pageref{LastPage}}
\renewcommand{\headrulewidth}{0.4pt}
\setlength{\parindent}{0pt}
\setlength{\parskip}{0.5em}
\setlist{itemsep=0.3em,topsep=0.3em}
\newcommand{\apartado}[1]{\subsection*{#1}}
\pgfplotsset{
  revision/.style={
    width=0.94\linewidth,height=7.5cm,
    axis lines=middle,grid=major,grid style={gray!18},
    tick label style={font=\small},label style={font=\small},
    xlabel={$x$},ylabel={$y$},samples=160,
    legend style={at={(0.5,-0.16)},anchor=north,draw=none,font=\small},
    clip=true
  }
}

\hypersetup{pdftitle={Matrices III - Soluciones revisadas}}
\begin{document}
\section*{Matrices III: soluciones}
Solucionario elaborado en esta revisión a partir de F016. El ZIP original no incluía soluciones.
\section*{Ejercicio 1(a): factor común y ecuación}
Como $X$ aparece a la derecha en ambos sumandos, la propiedad distributiva da
\[AX+BX=(A+B)X.\]
No se cambia el orden de los factores. Con las matrices del enunciado,
\[
H=A+B=\begin{pmatrix}1&0\\2&1\end{pmatrix},\qquad\det H=1.
\]
Existe inversa:
\[
H^{-1}=\begin{pmatrix}1&0\\-2&1\end{pmatrix},\qquad
HH^{-1}=H^{-1}H=I_2.
\]
Multiplicamos por ella a la izquierda:
\[
X=H^{-1}C=\begin{pmatrix}1&0\\-2&1\end{pmatrix}
\begin{pmatrix}2&1\\3&2\end{pmatrix}
=\boxed{\begin{pmatrix}2&1\\-1&0\end{pmatrix}}.
\]
\textbf{Comprobación en la ecuación original:}
\[
AX=\begin{pmatrix}1&0\\2&1\end{pmatrix},\quad
BX=\begin{pmatrix}1&1\\1&1\end{pmatrix},\quad
AX+BX=\begin{pmatrix}2&1\\3&2\end{pmatrix}=C.
\]
La solución es única porque $A+B$ es invertible.
\section*{Ejercicio 1(b): potencias}
\[
A^2=\begin{pmatrix}0&-1\\1&0\end{pmatrix}^2
=\begin{pmatrix}-1&0\\0&-1\end{pmatrix}=-I_2.
\]
Por tanto,
\[
\boxed{A^4=(A^2)^2=I_2},\qquad
\boxed{A^{80}=(A^4)^{20}=I_2}.
\]
Se conservan los exponentes 4 y 80 del original. El ciclo de potencias se repite cada cuatro exponentes.
\clearpage
\section*{Ejercicio 2: operaciones y dimensiones}
Se muestran las cinco respuestas para consulta. En el examen se deben elegir \textbf{solo cuatro}; el ejercicio completo vale 2,5 puntos.

Aquí $A,B$ son $2\times2$, $C$ es $2\times1$ y $D$ es $1\times2$. Se utilizan las matrices propias del ejercicio 2.
\subsection*{1. $AC^t+BD^t$}
$C^t$ tiene dimensión $1\times2$. El producto $AC^t$ no existe, porque las dimensiones internas son 2 y 1. Por ello la operación completa \textbf{no está definida}, aunque $BD^t$ sí exista.
\subsection*{2. $(A+B)C$}
\[
A+B=\begin{pmatrix}4&2\\-2&-1\end{pmatrix},\qquad
\boxed{(A+B)C=\begin{pmatrix}-10\\5\end{pmatrix}}.
\]
El producto tiene dimensión $2\times1$.
\subsection*{3. $B^2+CD$}
\[
B^2=\begin{pmatrix}10&0\\0&10\end{pmatrix},\qquad
CD=\begin{pmatrix}0&-6\\0&-3\end{pmatrix}.
\]
Ambas son $2\times2$:
\[\boxed{B^2+CD=\begin{pmatrix}10&-6\\0&7\end{pmatrix}}.\]
\subsection*{4. $|-2A|$}
Como $A$ es de orden 2,
\[|-2A|=(-2)^2|A|=4\bigl(2\cdot1-(-1)(-4)\bigr)=4(-2)=\boxed{-8}.\]
\subsection*{5. $|C^tC|$}
Este producto \textbf{sí está definido}: $(1\times2)(2\times1)$ da una matriz $1\times1$.
\[C^tC=\begin{pmatrix}-2&-1\end{pmatrix}\begin{pmatrix}-2\\-1\end{pmatrix}=\begin{pmatrix}5\end{pmatrix}.\]
El determinante de una matriz de orden 1 es su único elemento:
\[\boxed{|C^tC|=5}.\]
No se confunde con $|C|$, que no estaría definido por ser $C$ rectangular.
\clearpage
\section*{Ejercicio 3: igualdad $MA=A$}
La multiplicación a la izquierda por $M$ intercambia las dos filas:
\[
MA=\begin{pmatrix}0&1\\1&0\end{pmatrix}
\begin{pmatrix}a&b\\2&1\end{pmatrix}
=\begin{pmatrix}2&1\\a&b\end{pmatrix}.
\]
Imponemos la igualdad con la matriz original:
\[
\begin{pmatrix}2&1\\a&b\end{pmatrix}
=\begin{pmatrix}a&b\\2&1\end{pmatrix}.
\]
Igualando las entradas se obtiene $a=2$ y $b=1$; las condiciones de la segunda fila repiten las mismas ecuaciones.
\[\boxed{a=2,\qquad b=1}.\]
\textbf{Comprobación:}
\[
M\begin{pmatrix}2&1\\2&1\end{pmatrix}
=\begin{pmatrix}2&1\\2&1\end{pmatrix}=A.
\]
Las filas iguales permanecen invariantes al intercambiarlas. Esta matriz $A$ es singular, pero el ejercicio no exige que tenga inversa: la igualdad se verifica directamente por multiplicación.
\end{document}
